\begin{proof}
We use the extended-real convention stated before
\Cref{thm:ch16-collapse-finite-tropical-curves}.  The conditions defining
\(\mathcal J_X\) are local.  For the only nonliteral assertion, suppose the
junction inequality holds on an open cover and consider one admissible
tuple.  An open member containing its node contains, for all sufficiently
small \(s>0\), the segments from \((0,y_0)\) to
\[
 \bigl(st_av_{i(a)},\,y_0+s(y_a-y_0)\bigr).
\]
The local junction inequality at these contracted points, followed by
branchwise convexity, gives the asserted inequality after cancelling
\((1-s)u(0,y_0)\) and dividing by \(s\).  The case
\(u(0,y_0)=-\infty\) is automatic.  Thus \(\mathcal J_X\) is a sheaf.
It is closed under pointwise decreasing limits, since both extended-real
convexity and every finite junction inequality pass to such a limit.
Every generator belongs to \(\mathcal J_X\): its ambient affine tangent at
\((0,y_0)\) is a minorant on every branch, and its weighted linear terms
cancel at the displayed barycentre.  Hence
\[
                         \mathcal P_X\subseteq\mathcal J_X.
 \tag{\(\dagger\)}
\]

It remains to prove the reverse inclusion at the first stage.  Away from
\(\{0\}\times P\), this is the one-point case of
\Cref{thm:p071-stone-polyhedral-collapse}, with a singleton clopen
parameter and a sufficiently small convex polyhedral normal-coordinate
interval.  We work at \((0,y_0)\), with \(u\in\mathcal J_X\) defined on
a product neighbourhood.  Choose numbers \(0<2T_i<r_i\), bounded open
polyhedra \(O_1,O_0\subset A\), and put
\[
 V=P\cap O_1,\qquad Q=P\cap\overline{O_0},
\]
so that \(y_0\in V\), \(\overline V^{\,P}\subset Q\), \(Q\) is a
compact convex polytope, and
\[
 D:=\Bigl(\{0\}\cup\coprod_i(0,2T_i]v_i\Bigr)\times Q
\]
is contained in the original domain.  We first construct continuous finite
comparison majorants of \(u|_D\).  This is the step which treats every
polyhedral face and corner.

For a continuous real-valued function \(\phi\) on \(D\), define its
\emph{junction-convex envelope} \(C\phi\) as follows.  At a point
\(x=(tv_i,y)\) with \(t>0\), take the infimum of
\[
                   \sum_a\alpha_a\phi(s_av_{j(a)},y_a)
 \tag{\(\clubsuit\)}
\]
over all finite tuples in \(D\) satisfying
\[
 \begin{aligned}
 \sum_a\alpha_a&=1,
 &\sum_a\alpha_as_av_{j(a)}&=tv_i,\\
 \sum_a\alpha_ay_a&=y,
 &\sum_{j(a)=i}\alpha_as_a&\geq t.
 \end{aligned}
 \tag{\(\clubsuit_1\)}
\]
At a node \((0,y)\), use the same definition but omit the last inequality;
it is then independent of a branch label.  The singleton tuple shows that
\(C\phi\leq\phi\).

We record the three properties of this envelope needed below.  First,
\(C\phi\) satisfies the two comparison conditions on \(D\).  Indeed,
convex combinations of admissible tuples prove branchwise convexity, and
combining admissible tuples for a balanced node family proves the junction
inequality.  Second, it is the largest function satisfying those conditions
below \(\phi\).  To see this, let \(v\) satisfy them with \(v\leq\phi\),
and consider a tuple in
\((\clubsuit_1)\).  Put
\[
               \lambda=\sum_{j(a)=i}\alpha_as_a\geq t.
\]
For \(t>0\), retain the proportion \(t/\lambda\) of each branch-\(i\)
weight and call all remaining weights residual.  The retained weights have
normal moment \(tv_i\), while the residual weights have normal moment zero.
If the residual mass is nonzero, apply the junction inequality to its
Euclidean barycentre in \(Q\), and then apply branchwise convexity to that
node and the retained branch-\(i\) points.  This gives
\[
                  v(tv_i,y)\leq
                  \sum_a\alpha_av(s_av_{j(a)},y_a).
 \tag{\(\clubsuit_2\)}
\]
If the residual mass is zero, this is ordinary branchwise convexity.  At a
node, \((\clubsuit_2)\) is precisely the junction inequality.  Taking the
infimum proves \(v\leq C\phi\).

Third, \(C\phi\) is finite and continuous on \(D\).  For attainment and
lower semicontinuity, regard the tuples as probability measures on the
compact space \(D\).  Their Euclidean and normal moment constraints are
closed.  The extra branch-\(i\) moment is the integral of the continuous
function which equals \(s\) on branch \(i\), zero on the other branches,
and zero at the common node.  Compactness therefore supplies a minimizing
measure and shows lower semicontinuity under convergent target points; at a
node pass to one of the finitely many branch labels.  Carath\'eodory applied
also to the cost coordinate replaces every minimizing measure by a finite
one, as used in \((\clubsuit)\).  The function \(C\phi\) is finite and
convex on each compact branch prism \([0,2T_i]\times Q\).  A finite convex
function on a compact polytope is upper semicontinuous: triangulate near a
point with that point as a vertex and use the barycentric coefficient tending
to one.  Thus lower and upper semicontinuity give continuity on each branch,
and the common node values glue the finitely many branches continuously.

The restriction to \(D\) of every member of \(\mathcal J_X\) is upper
semicontinuous.  Indeed, \(\max\{u,-N\}\) is finite and satisfies the two
comparison conditions on \(D\): the maximum operation is convex and
nondecreasing, so Jensen's inequality preserves both defining conditions.
It is upper semicontinuous on each compact branch prism by the preceding
polyhedral argument, hence on \(D\); taking the decreasing infimum over
\(N\) proves the assertion for \(u\).

Fix a compatible metric \(d\) on \(D\).  The standard compact-metric
upper regularization
\[
 \phi_n(x):=\sup_{z\in D}
 \left(\max\{u(z),-n\}-n\,d(x,z)\right)
 \tag{\(\clubsuit_3\)}
\]
is finite, continuous, decreases to \(u\), and dominates \(u\).  For the
last assertion, an upper bound for \(u\) on \(D\) and upper
semicontinuity separate a small ball about the chosen point from its compact
complement; the penalty in \((\clubsuit_3)\) then forces convergence.  Set
\[
                              H_n:=C\phi_n.
\]
The preceding three properties give
\[
       u\leq H_n\leq\phi_n,\qquad H_{n+1}\leq H_n,\qquad H_n\downarrow u,
 \tag{\(\clubsuit_4\)}
\]
and every \(H_n\) is finite, continuous, and satisfies the two comparison
conditions on \(D\).

We next give a uniform smooth majorant for one such continuous comparison
function \(H:=H_n\).  Write \(H_i(t,y)=H(tv_i,y)\), including
\(H_i(0,y)=H(0,y)\).  After choosing an origin in \(A\), for
\(p\in A^*\) put
\[
 f_i(t,p)=\min_{y\in Q}\bigl(H_i(t,y)-p\mathbin\cdot y\bigr),
 \qquad
 f_0(p)=\min_{y\in Q}\bigl(H(0,y)-p\mathbin\cdot y\bigr).
 \tag{\(\spadesuit\)}
\]
These are finite continuous scalar star profiles, convex in \(t\), with
common central value \(f_0(p)\).  The product junction inequality proves
their finite-star junction inequalities exactly as follows: for a balanced
normal tuple choose \(y_a\in Q\), put \(\bar y=\sum_a\alpha_ay_a\), and
write
\[
 \begin{split}
 f_0(p)
 &\leq H(0,\bar y)-p\mathbin\cdot\bar y\\
 &\leq\sum_a\alpha_a
      \bigl(H_{i(a)}(t_a,y_a)-p\mathbin\cdot y_a\bigr).
 \end{split}
 \tag{\(\spadesuit_1\)}
\]
Taking independent infima proves the claim.

For every fixed \(t\), the function obtained from \(H_i(t,\cdot)\) by
setting it equal to \(+\infty\) off \(Q\) is proper, lower semicontinuous,
and convex: continuity on the closed polytope is exactly what is needed at
its faces.  Fenchel--Moreau therefore gives, including every boundary point,
\[
 H_i(t,y)=\sup_{p\in A^*}
               \bigl(p\mathbin\cdot y+f_i(t,p)\bigr)
 \qquad (y\in Q).
 \tag{\(\spadesuit_2\)}
\]
Choose nested finite partitions of \([0,2T_i]\), independently of \(p\),
with mesh tending to zero.  Their scalar chordal interpolants
\(h_{\ell,i}(t,p)\) majorize \(f_i(t,p)\), decrease under refinement, and
satisfy the finite-star junction inequalities.  Let
\[
 \widehat H_{\ell,i}(t,y)=
 \sup_p\bigl(p\mathbin\cdot y+h_{\ell,i}(t,p)\bigr).
 \tag{\(\spadesuit_3\)}
\]
The same calculation as in \((\spadesuit_1)\) shows that
\(\widehat H_\ell\) satisfies the two comparison conditions on \(D\).  If
\(t\) lies between adjacent
partition points \(a,b\), with affine weights \(1-\theta,\theta\), then
\[
 H_i(t,y)\leq\widehat H_{\ell,i}(t,y)
 \leq(1-\theta)H_i(a,y)+\theta H_i(b,y).
 \tag{\(\spadesuit_4\)}
\]
Uniform continuity of \(H\) on \(D\) therefore permits a choice of
\(\ell\) for which
\[
                       H\leq\widehat H_\ell\leq H+\varepsilon
 \quad\text{on }D.
 \tag{\(\spadesuit_5\)}
\]
The function \(\widehat H_\ell\) is lower semicontinuous as a supremum and
finite convex on every compact branch prism; the preceding polyhedral
argument makes it continuous on \(D\).

Compactness now yields covectors \(p_1,\ldots,p_q\) such that, with the
middle expression read on branch \(i\),
\[
 \widehat H_\ell-\varepsilon
 \leq A:=\max_j\bigl(p_j\mathbin\cdot y+h_{\ell,i}(t,p_j)\bigr)
 \leq\widehat H_\ell
 \quad\text{on }D.
 \tag{\(\spadesuit_6\)}
\]
For each of the finitely many scalar profiles, the quantitative
common-affine-minorant and regularized-maximum construction in the proof of
\Cref{thm:ch16-collapse-finite-tropical-curves} supplies a smooth ambient
branchwise-convex majorant \(g_j\) on the fixed smaller substar with
\[
                  h_\ell(\cdot,p_j)\leq g_j
                  \leq h_\ell(\cdot,p_j)+\varepsilon.
\]
The functions \(g_j(tv_i)+p_j\mathbin\cdot y\) are ambient-smooth and
jointly convex in \((t,y)\).  Taking their convex coordinatewise
nondecreasing regularized maximum, with error \(\varepsilon\), and adding
one further \(\varepsilon\), gives a function
\(w\in\mathcal G_{X,0}\) on the fixed product
\[
 V_X:=\Bigl(\{0\}\cup\coprod_i(0,\rho_i)v_i\Bigr)\times V,
 \qquad 0<\rho_i<T_i,
\]
such that
\[
                         H\leq w\leq H+4\varepsilon
 \quad\text{on }V_X.
 \tag{\(\spadesuit_7\)}
\]
The affine collar at the common node exists simultaneously because only
finitely many scalar profiles occur; it may depend on \(n\), while the
displayed product neighbourhood does not.  All functions are ambient smooth
in the \(A\)-variables, so this step also applies on every face of \(P\).

Apply this construction to \(H_n\) with a summable positive sequence
\((\varepsilon_n)\), producing \(w_n\) on the one fixed \(V_X\).  Put
\[
 d_n:=4\sum_{k>n}\varepsilon_k,
 \qquad G_n:=w_n+d_n.
\]
Equations \((\clubsuit_4)\) and \((\spadesuit_7)\) give
\[
 G_{n+1}\leq H_{n+1}+4\varepsilon_{n+1}+d_{n+1}
 \leq H_n+d_n\leq G_n.
\]
Since \(H_n\downarrow u\) and \(4\varepsilon_n+d_n\to0\), the sequence
\(G_n\) decreases pointwise to \(u\) on \(V_X\).  Thus
\(\mathcal J_X\subseteq\mathcal P_{X,1}\).  Combining this with
\((\dagger)\) and \(\mathcal P_{X,1}\subseteq\mathcal P_X\) proves the
theorem.
\end{proof}
